% 精简版 chapters.tex — 基于考试重点筛选
% 生成命令: uv run python scripts/generate_filtered.py
% 总页数: 439


% ====== 引言（共 9 页） ======
% \input{notes/1引言-notes.tex}  % <-- 章末笔记

\slidegrid{1引言}{5}{39}{40}{41}{42}{43}{44}{45}
\slidegrid{1引言}{46}{0}{0}{0}{0}{0}{0}{0}

% ── 引言 章末笔记 ──


% ====== 数学基础（共 13 页） ======
% \input{notes/2数学基础-notes.tex}  % <-- 章末笔记

\slidegrid{2数学基础}{3}{4}{5}{6}{7}{8}{9}{10}
\slidegrid{2数学基础}{11}{12}{13}{14}{15}{0}{0}{0}

% ── 数学基础 章末笔记 ──


% ====== 凸分析（共 28 页） ======
% \input{notes/3凸分析-notes.tex}  % <-- 章末笔记

\slidegrid{3凸分析}{3}{4}{5}{8}{9}{11}{12}{14}
\slidegrid{3凸分析}{15}{17}{18}{21}{23}{24}{25}{27}
\slidegrid{3凸分析}{29}{30}{34}{39}{40}{41}{42}{43}
\slidegrid{3凸分析}{44}{45}{46}{47}{0}{0}{0}{0}

% ── 凸分析 章末笔记 ──


% ====== 线性规划基本性质（共 26 页） ======
% \input{notes/4线性规划基本性质-notes.tex}  % <-- 章末笔记

\slidegrid{4线性规划基本性质}{3}{4}{5}{6}{7}{8}{9}{10}
\slidegrid{4线性规划基本性质}{11}{12}{13}{14}{15}{16}{17}{21}
\slidegrid{4线性规划基本性质}{22}{23}{24}{25}{26}{27}{28}{29}
\slidegrid{4线性规划基本性质}{30}{31}{0}{0}{0}{0}{0}{0}

% ── 线性规划基本性质 章末笔记 ──


% ====== 单纯形方法（共 40 页） ======
% \input{notes/5单纯形方法-notes.tex}  % <-- 章末笔记

\slidegrid{5单纯形方法}{2}{3}{4}{5}{6}{7}{12}{15}
\slidegrid{5单纯形方法}{16}{24}{25}{26}{27}{28}{29}{30}
\slidegrid{5单纯形方法}{35}{36}{37}{38}{39}{40}{41}{42}
\slidegrid{5单纯形方法}{43}{44}{45}{46}{47}{48}{49}{50}
\slidegrid{5单纯形方法}{51}{52}{53}{54}{55}{56}{57}{58}

% ── 单纯形方法 章末笔记 ──


% ====== 对偶理论（共 57 页） ======
% \input{notes/6对偶理论-notes.tex}  % <-- 章末笔记

\slidegrid{6对偶理论}{3}{4}{5}{6}{7}{8}{9}{10}
\slidegrid{6对偶理论}{11}{15}{16}{17}{18}{19}{21}{22}
\slidegrid{6对偶理论}{23}{24}{25}{26}{27}{28}{29}{30}
\slidegrid{6对偶理论}{31}{32}{33}{40}{41}{42}{43}{44}
\slidegrid{6对偶理论}{45}{46}{47}{48}{49}{50}{51}{52}
\slidegrid{6对偶理论}{53}{54}{55}{56}{57}{58}{59}{60}
\slidegrid{6对偶理论}{61}{62}{63}{64}{65}{66}{67}{68}
\slidegrid{6对偶理论}{81}{0}{0}{0}{0}{0}{0}{0}

% ── 对偶理论 章末笔记 ──


% ====== 最优性条件（共 43 页） ======
% \input{notes/7最优性条件-notes.tex}  % <-- 章末笔记

\slidegrid{7最优性条件}{3}{4}{5}{7}{13}{14}{15}{16}
\slidegrid{7最优性条件}{17}{18}{19}{20}{21}{22}{23}{24}
\slidegrid{7最优性条件}{25}{26}{27}{28}{29}{34}{35}{36}
\slidegrid{7最优性条件}{37}{38}{39}{40}{41}{42}{43}{48}
\slidegrid{7最优性条件}{49}{50}{51}{52}{53}{54}{55}{56}
\slidegrid{7最优性条件}{57}{58}{59}{0}{0}{0}{0}{0}

% ── 最优性条件 章末笔记 ──

\notefullpage{
  \section*{FJ条件 vs KKT条件}

  \subsection*{Fritz John 条件}
  存在不全为零的 $\lambda_0,\lambda_1,\dots,\lambda_m,\mu_1,\dots,\mu_l$ 使得
  \begin{align*}
    \lambda_0 \nabla f(x^*) + \sum_{i=1}^{m} \lambda_i \nabla g_i(x^*)
    + \sum_{j=1}^{l} \mu_j \nabla h_j(x^*) &= 0 \\
    \lambda_i &\ge 0,\quad i=0,1,\dots,m \\
    \lambda_i g_i(x^*) &= 0,\quad i=1,\dots,m
  \end{align*}
  \textbf{注意：}$\lambda_0$ 可以为 0！此时不含目标函数信息，FJ点可能无意义。

  \subsection*{KKT 条件（加约束规格后）}
  存在 $\lambda_1,\dots,\lambda_m,\mu_1,\dots,\mu_l$ 使得
  \begin{align*}
    \nabla f(x^*) + \sum_{i=1}^{m} \lambda_i \nabla g_i(x^*)
    + \sum_{j=1}^{l} \mu_j \nabla h_j(x^*) &= 0 \\
    \lambda_i &\ge 0,\quad i=1,\dots,m \\
    \lambda_i g_i(x^*) &= 0,\quad i=1,\dots,m
  \end{align*}
  即 Lagrange 函数 $L(x,\lambda,\mu) = f(x) + \sum\lambda_i g_i(x) + \sum\mu_j h_j(x)$ 的驻点条件。

  \textbf{核心区别：}FJ中$\lambda_0$可能为0（不含目标函数），KKT通过约束规格保证$\lambda_0 \neq 0$。

  真题链接：2025 Q6 —— 先求所有FJ点，再证不存在KKT点。
}

% ====== 算法（共 19 页） ======
% \input{notes/8算法-notes.tex}  % <-- 章末笔记

\slidegrid{8算法}{3}{4}{5}{6}{7}{8}{9}{10}
\slidegrid{8算法}{14}{15}{16}{17}{18}{19}{20}{21}
\slidegrid{8算法}{22}{23}{24}{0}{0}{0}{0}{0}

% ── 算法 章末笔记 ──


% ====== 一维搜索（共 17 页） ======
% \input{notes/9一维搜索-notes.tex}  % <-- 章末笔记

\slidegrid{9一维搜索}{3}{4}{5}{6}{7}{26}{27}{51}
\slidegrid{9一维搜索}{52}{53}{54}{55}{56}{57}{58}{59}
\slidegrid{9一维搜索}{60}{0}{0}{0}{0}{0}{0}{0}

% ── 一维搜索 章末笔记 ──


% ====== 使用导数的优化算法（共 79 页） ======
% \input{notes/10使用导数的优化算法-notes.tex}  % <-- 章末笔记

\slidegrid{10使用导数的优化算法}{40}{41}{42}{43}{44}{45}{46}{47}
\slidegrid{10使用导数的优化算法}{48}{49}{50}{51}{52}{53}{54}{55}
\slidegrid{10使用导数的优化算法}{56}{57}{58}{59}{60}{61}{62}{63}
\slidegrid{10使用导数的优化算法}{64}{65}{66}{67}{68}{69}{70}{71}
\slidegrid{10使用导数的优化算法}{72}{73}{74}{75}{76}{77}{78}{79}
\slidegrid{10使用导数的优化算法}{80}{81}{82}{83}{84}{85}{86}{87}
\slidegrid{10使用导数的优化算法}{88}{89}{90}{91}{92}{93}{94}{95}
\slidegrid{10使用导数的优化算法}{96}{97}{98}{99}{100}{101}{102}{103}
\slidegrid{10使用导数的优化算法}{104}{105}{106}{107}{108}{109}{110}{111}
\slidegrid{10使用导数的优化算法}{112}{113}{114}{115}{116}{117}{118}{0}

% ── 使用导数的优化算法 章末笔记 ──

\notefullpage{
  \section*{拟牛顿法公式速查}

  \subsection*{DFP法（对Hessian逆矩阵$H$的秩2校正）}
  \begin{align*}
    s_k &= x_{k+1} - x_k,\qquad y_k = g_{k+1} - g_k \\
    H_{k+1} &= H_k + \frac{s_k s_k^T}{s_k^T y_k}
              - \frac{H_k y_k y_k^T H_k}{y_k^T H_k y_k}
  \end{align*}
  搜索方向: $d_k = -H_k \nabla f(x_k)$\\[2pt]
  算法步骤: ①给定$x_0$, $H_0=I$ ②计算$g_0$, $d_0=-H_0g_0$ ③一维搜索求$\lambda_0$ ④$x_1=x_0+\lambda_0 d_0$ ⑤计算$s_0,y_0$ ⑥DFP校得$H_1$ ⑦重复。

  \subsection*{BFGS法（对Hessian矩阵$B$的秩2校正）}
  \begin{align*}
    B_{k+1} &= B_k + \frac{y_k y_k^T}{y_k^T s_k}
              - \frac{B_k s_k s_k^T B_k}{s_k^T B_k s_k}
  \end{align*}
  搜索方向: 解$B_k d_k = -\nabla f(x_k)$ （或直接用$H_k$的BFGS公式）\\[2pt]
  BFGS也可以写成对$H$的校正公式（镜像形式）。

  \textbf{关键区别：} DFP校$H$（Hessian逆），BFGS校$B$（Hessian本身）。

  真题链接：2025 Q4 —— DFP法完整两轮迭代。
}

% ====== 无约束优化直接方法（共 0 页） ======
% \input{notes/11无约束优化直接方法-notes.tex}  % <-- 章末笔记

% （本章不考，已全部排除）


% ====== 可行方向法（共 60 页） ======
% \input{notes/12可行方向法-notes.tex}  % <-- 章末笔记

\slidegrid{12可行方向法}{2}{3}{4}{5}{6}{7}{8}{9}
\slidegrid{12可行方向法}{10}{11}{12}{13}{14}{15}{16}{17}
\slidegrid{12可行方向法}{18}{19}{20}{21}{22}{23}{24}{25}
\slidegrid{12可行方向法}{26}{27}{28}{29}{30}{31}{32}{33}
\slidegrid{12可行方向法}{39}{40}{41}{42}{43}{44}{45}{46}
\slidegrid{12可行方向法}{47}{48}{49}{50}{51}{52}{53}{54}
\slidegrid{12可行方向法}{55}{56}{57}{58}{59}{60}{61}{62}
\slidegrid{12可行方向法}{63}{64}{65}{66}{0}{0}{0}{0}

% ── 可行方向法 章末笔记 ──

\notefullpage{
  \section*{Rosen梯度投影法 算法步骤}

  \subsection*{核心公式}
  投影矩阵: $P = I - M^T(MM^T)^{-1}M$
  搜索方向: $d = -P\,\nabla f(x)$

  \subsection*{算法流程}
  \begin{enumerate}
    \item 在当前点$x_k$识别\textbf{积极约束}（取等号的约束），构造积极约束系数矩阵$M$
    \item 计算投影矩阵 $P = I - M^T(MM^T)^{-1}M$
    \item 计算搜索方向 $d_k = -P \nabla f(x_k)$
    \item \textbf{若$d_k \neq 0$}：沿$d_k$做一维搜索，求步长$\lambda_k$，
          $x_{k+1} = x_k + \lambda_k d_k$，回到步骤1
    \item \textbf{若$d_k = 0$}：计算$w = -(MM^T)^{-1}M\nabla f(x_k)$
          \begin{itemize}
            \item 若$w$的所有分量$\ge 0$ → $x_k$为\textbf{KKT点}，停止
            \item 若$w$有负分量 → 去掉对应的积极约束行，重新计算$P$
          \end{itemize}
  \end{enumerate}

  \subsection*{Zoutendijk法（补充）}
  解LP子问题求下降可行方向：
  \[
  \begin{aligned}
  \min\ &\nabla f(x)^T d \\
  \text{s.t.}\ &\nabla g_i(x)^T d + g_i(x) \le 0,\quad i\in I(x) \\
  &-1 \le d_j \le 1
  \end{aligned}
  \]
  若最优值$<0$得下降可行方向；若$=0$则为FJ/KKT点。

  真题链接：2025 Q5 —— Rosen梯度投影法求解约束优化。
}

% ====== 罚函数法（共 48 页） ======
% \input{notes/13罚函数法-notes.tex}  % <-- 章末笔记

\slidegrid{13罚函数法}{2}{3}{4}{5}{6}{7}{8}{9}
\slidegrid{13罚函数法}{10}{11}{12}{13}{14}{15}{16}{17}
\slidegrid{13罚函数法}{18}{19}{20}{21}{22}{23}{24}{25}
\slidegrid{13罚函数法}{26}{27}{28}{29}{30}{31}{32}{33}
\slidegrid{13罚函数法}{34}{35}{36}{37}{38}{39}{40}{41}
\slidegrid{13罚函数法}{42}{43}{44}{45}{46}{47}{48}{49}

% ── 罚函数法 章末笔记 ──

